\section{Introduction}
\label{sec:introduction}

Large language models (LLMs) are widely deployed in production services such as conversational assistants~\cite{XX}, search augmentation~\cite{XX}, code generation~\cite{XX}, and domain-specific intelligent applications~\cite{XX}. As a result, LLM serving has emerged as one of the core GPU workloads in modern datacenters. Serving LLMs at scale consumes substantial GPU resources, creates significant energy and cost pressure, and must still satisfy strict user-facing latency targets such as time to first token (TTFT) and time per output token (TPOT). Modern GPUs support dynamic voltage and frequency scaling (DVFS), offering a natural hardware mechanism for improving energy efficiency. However, applying DVFS effectively to LLM serving under strict latency constraints remains challenging. 

To reduce energy consumption under strict latency constraints, prior work has explored increasingly fine-grained DVFS control for LLM serving.
%
DynamoLLM~\cite{XX} manages energy at the service level by jointly tuning instance count, model parallelism, and GPU frequency, while throttLL’eM~\cite{XX} performs predictive GPU throttling based on iteration-level serving dynamics such as batch size and KV-cache usage. BiScale~\cite{XX} further exploits the fundamental distinction between prefill and decode by assigning separate frequency settings to the two phases. These studies show that DVFS can reduce the energy consumption of LLM serving under latency constraints, and phase-aware designs further demonstrate the value of exploiting the distinct characteristics of prefill and decode. However, phase-level control remains too coarse, as it overlooks operator-level heterogeneity within each phase.

Attention and feed-forward networks (FFNs), the two dominant operator classes in LLM inference, place different demands on compute and memory resources, and these demands vary across serving phases, workloads, and system configurations. In prefill, both operators can shift between memory-bound and compute-bound behavior depending on tensor parallelism, batch size, and related workload parameters. In decode, Attention becomes increasingly memory-bound as tensor parallelism and batch size increase because of KV-cache accesses and declining arithmetic intensity, whereas FFNs remain compute-sensitive over a much wider range of configurations due to their reliance on dense matrix multiplications. These bottleneck shifts directly affect DVFS behavior: increasing compute frequency benefits compute-bound operators more, whereas memory-bound operators remain constrained primarily by memory access behavior. Consequently, a single DVFS setting cannot adequately serve both Attention and FFNs across phases, or even within the same phase under different workloads and system configurations.

However, exploiting this operator-level heterogeneity through Attention/FFN (A/F) disaggregation and A/F-aware DVFS is nontrivial.
%
First, prior work~\cite{XX} has shown that Attention and FFNs disaggregation (AFD) incurs additional inter-operator data transfers at every layer boundary, which can reduce throughput if not carefully managed.
%
Second, AFD raises a heterogeneous resource provisioning challenge, since Attention and FFNs stress memory and compute resources differently and may therefore require different resource allocations and coordination.
%
Third, AFD requires coordinated frequency control across Attention and FFNs, as their frequencies must be jointly selected with transfer overheads to minimize energy while satisfying strict latency SLOs.
%
Finally, flexible A/F-aware DVFS requires switching-aware runtime control, as non-negligible clock-switching overheads can otherwise offset its energy-efficiency gains.

\cc{need re-writing}

% \cc{To address these challenges, we present AFlex, an operator-aware, SLO-aware serving framework for LLM inference. AFlex introduces a low-overhead cross-operator state-transfer path to reduce the transfer overhead of Attention--FFN (A/F) disaggregation, a heterogeneous provisioning mechanism to accommodate the distinct resource demands of Attention and FFNs, a coordinated frequency selection algorithm to jointly optimize their frequencies under strict latency SLOs, and a switching-aware runtime scheduler to account for non-negligible clock-switching overheads on real GPUs. AFlex also includes Lex, an optimized unified-frequency selection algorithm that serves as a strong single-frequency baseline and helps isolate the gains of finer-grained control. Compared with Lex, AFlex achieves an additional 5\%--14\% energy reduction and up to 39\% end-to-end energy savings under measured clock-switching overheads on real GPUs, where the median switching latency is 4.5\,ms.


% This paper makes four contributions. First, we identify and characterize operator-level frequency heterogeneity between Attention and FFNs in LLM serving. Second, we develop Lex, an optimized unified-frequency selection algorithm that strengthens prior phase-level DVFS and serves as a strong unified baseline. Third, we design and implement AFlex, an operator-aware, SLO-aware serving framework that extends Lex with A/F disaggregation, a low-overhead cross-operator state-transfer path, and a switching-aware scheduler. Fourth, we evaluate Lex and AFlex on Qwen3-32B running on A800 GPUs and show that unified DVFS remains suboptimal even under optimized single-frequency control, while AFlex delivers substantial energy savings across practical serving regimes. We next review the background on LLM serving, GPU DVFS, and phase-disaggregated serving.}